\section{Ordered arrays and conditional locality}\label{sec:arrays}

The next construction retains the order in which limits are taken.  Its
basic independence is along the last axis.  Additional independence and
conditional locality will be consequences of an extension argument; they
are not symmetry assumptions on the arrays.

\subsection{The initial array laws}

We interpret $I^0$ as a singleton.  For each $n\geq0$ there is a
probability law $\rho_n$ on $X^{I^n}$ with the following properties:
\begin{enumerate}
\item $\rho_0=\mu_X$, and $\rho_n$ is invariant under simultaneous $T_X$;
\item fixing any one axis at any slot gives $\rho_{n-1}$;
\item any proper family of last-axis planes is independent under
      $\rho_n$;
\item
\begin{equation}\label{eq:array-witness}
 \int_{X^I}\prod_{i\in I}f_i(x_i)\dd\rho_1(x)\geq\delta.
\end{equation}
\end{enumerate}
Here a last-axis plane is the whole array with its last coordinate fixed;
``proper'' means containing at most $k-1$ of the $k$ planes.

To construct the laws, fix a free ultrafilter on the layout indices $b$.
For $b_1,\ldots,b_n$, start with the law of
\[
 \left(T_X^{\sum_{\ell=1}^n d_{j_\ell}(b_\ell)}x
       \right)_{(j_1,\ldots,j_n)\in I^n},\qquad x\sim\mu_X.
\]
Take the ultralimit in $b_n$ first, then $b_{n-1}$, and so on.  These
limits exist on cylinder probabilities.  They are consistent
probabilities on every finite collection of alphabet coordinates and
therefore determine a probability measure on the countable product
$X^{I^n}$.  These are also successive weak ultralimits on this compact
metrizable space.  Ultralimits of bounded scalar sequences are
linear and multiplicative.

Time invariance is immediate.  Fixing one axis adds a common time
translation to all remaining entries; invariance of $\mu_X$ removes that
translation before the limits are taken.  This proves consistency for
every axis.  To verify last-axis independence, take cylinder tests on
$s\leq k-1$ distinct last-axis planes and hold $b_1,\ldots,b_{n-1}$
fixed.  Each test is then a fixed bounded function of $x$, translated by
the corresponding $d_j(b_n)$.  Mixing of order $s$ makes their product
integral converge to the product of their integrals.  Multiplicativity
of the remaining ultralimits gives independence with plane marginal
$\rho_{n-1}$; for $s=1$ this is simply the marginal assertion and
requires no mixing hypothesis. Cylinder tests determine the asserted law. Finally,
\eqref{eq:array-witness} follows from \eqref{eq:witness-correlation}.

\subsection{The class of extensions}

All spaces in the rest of the proof are standard Borel spaces, with
their probability completions when discussing measurable functions or
$\sigma$-fields.  Conditional laws are regular conditional laws.  Their
identities are equalities almost surely, determined first on countable
generating families of tests.  This permits simultaneous choices of
versions for the countably many array levels and identities below.

At level $1$, the law $\rho_1$ still carries
\eqref{eq:array-witness}. At level $2$, view $I^2$ as a square: proper
families of its columns, the last-axis planes, are independent as
whole column arrays. We have not asserted the analogous row
independence. The higher levels will let us roll a first-axis row
into one local state and retain the same last-axis contract. To do
that we enlarge every $\rho_n$ compatibly, retaining an $X$-factor and
recording a parameter at each local state.

\begin{definition}[An array object]\label{def:array-object}
An \emph{array object} consists of standard probability systems
$(Y,m_0,T_Y)$ and $(U,\nu,T_U)$, factor maps
\[
 \kappa:Y\longrightarrow X,\qquad u:Y\longrightarrow U,
\]
and probability laws $m_n$ on $Y^{I^n}$ for $n\geq0$.  All
transformations are invertible and the factor maps commute with time.
The laws satisfy the following requirements.
\begin{enumerate}
\item Each $m_n$ is invariant under simultaneous $T_Y$ and projects
coordinatewise to $\rho_n$.  For $n\geq1$, fixing any axis at a slot
gives $m_{n-1}$, and every proper family of last-axis planes is
independent under $m_n$.
\item For $n\geq1$, conditional on the complete parameter array
$(u(y_v))_{v\in I^n}$, every proper family of last-axis planes is
independent.  The conditional law of each plane is $m_{n-1}$
conditioned only on that plane's own parameter array.
\item Under $m_0$, $\kappa(y)$ and $u(y)$ are independent; in particular,
$\kappa_*m_0=\mu_X$.
\end{enumerate}
An \emph{extension} is another such object together with a local factor
map $\theta:Y'\to Y$ preserving every array law and the map to $X$,
such that the old parameter is a factor of the new parameter.  The latter
means that $u\circ\theta=\psi\circ u'$ for a time-commuting factor
$\psi:U'\to U$.
\end{definition}

We also write $x(y)=\kappa(y)$ for the $X$-coordinate of a local state.
The initial laws give an object with $Y=X$ and a constant parameter.
For an object and $n\geq1$, let
\[
 \mathcal P_n=\sigma(u(y_v):v\in I^n),\qquad
 \mathcal O_{n,j}=\sigma(y_v:v_n\ne j),\qquad
 \mathcal G_{n,j}=\mathcal P_n\vee\mathcal O_{n,j}.
\]
Thus $\mathcal P_n$ records every parameter, $\mathcal O_{n,j}$ records
every state outside plane $j$, and $\mathcal G_{n,j}$ combines both.
Saturation will prevent an extension from revealing more about an old
test through any of these three kinds of data.

The construction uses the conditional-expectation energy and inverse-limit
method of sated extensions; see
\cite[preprint version, proof of Theorem~3.11 and Lemma~3.12]{Austin2015Sated}.
We verify directly that the defining conditional product laws of our
array category survive inverse limits.

\begin{proposition}[Saturation]\label{prop:saturation}
Every array object has an extension with the following property.  For
every $n\geq1$, every $F\in L^2(m_n)$, and each
\[
 \mathcal D\in\{\mathcal P_n\}\cup
 \{\mathcal G_{n,j},\mathcal O_{n,j}:j\in I\},
\]
no extension increases the norm of $\E[F\mid\mathcal D]$ after
pullback.  More precisely, if $\theta_n$ is the coordinatewise factor
map from any extension and $\mathcal D'$ is the corresponding data in
that extension, then
\begin{equation}\label{eq:saturation-equality}
 \E[F\circ\theta_n\mid\mathcal D']
 =\E[F\mid\mathcal D]\circ\theta_n.
\end{equation}
\end{proposition}

\begin{proof}
We first check closure under sequential inverse limits.  For a sequence
of extensions, take the local state and parameter spaces to be the
spaces of compatible paths, with their consistent probability laws.
They are standard Borel spaces.  At every array level the consistent
finite-stage laws determine a law on compatible array paths.  Time acts
coordinatewise.  Projection to $X$, consistency under fixing any axis,
and unconditional last-axis independence pass to the limit by testing
finite-stage functions.  Independence of the local $X$-state and the
limiting parameter follows in the same way.

For the conditional requirement, fix a proper family $J$ of last-axis
planes and bounded functions $F_j$ on those planes, all pulled from one
finite stage.  Let $\mathcal P_n^{(s)}$ denote the full parameter array
at stage $s$, and $\mathcal P_{n-1,j}^{(s)}$ its restriction to plane $j$.
For every sufficiently large $s$, the defining property gives
\[
 \E\left[\prod_{j\in J}F_j\,\middle|\,\mathcal P_n^{(s)}\right]
 =\prod_{j\in J}\E[F_j\mid\mathcal P_{n-1,j}^{(s)}].
\]
Increasing martingale convergence on both sides gives the same identity
for the limiting parameter fields.  The factors are uniformly bounded,
so their products converge in $L^2$.  This proves both independence
and dependence of each plane law only on its own parameters, first for
finite-stage tests and then, by density, for all bounded tests.

For a fixed test and data type, the norm of the indicated conditional
expectation is nondecreasing under extension and bounded above by the
norm of the test.  At each constructed stage, choose countable dense
families of bounded tests in every $L^2(m_n)$, and list all pairs of a
test and one of the finitely many data types at that level.  Construct
successive extensions as follows.  At round $q$, process the first $q$
listed pairs from each of the first $q$ stages already constructed.
For each pair choose an extension whose projection norm is within
$2^{-q}$ of the supremum over extensions of the current stage.
There is no requirement to attain the supremum.  Append these finitely
many extensions and continue.  Every listed pair is processed
infinitely often with error tending to zero.  Taking the sequential
inverse limit gives an object $M_\infty$ in the class.

Let $Z$ be any extension of $M_\infty$, and fix a test listed at some
finite stage.  At any sufficiently late processing step, $Z$ is also an
extension of the current stage.  Its projection norm is consequently
at most that step's supremum, which is at most the chosen next-stage
norm plus $2^{-q}$.  The next-stage norm is at most the limiting norm.
Letting $q\to\infty$ proves that $Z$ cannot increase the norm for this
test.  Finite-stage functions are dense in the inverse-limit
$L^2$ spaces, and conditional expectation is a contraction.  The same
conclusion follows for every $F\in L^2(m_n)$ of $M_\infty$.

Finally, the pullback of $\mathcal D$ is contained in $\mathcal D'$.
The Pythagorean identity for these nested orthogonal projections is
\begin{align*}
 &\norm{\E[F\circ\theta_n\mid\mathcal D']}_2^2
 -\norm{\E[F\mid\mathcal D]\circ\theta_n}_2^2\\
 &\hspace{2cm}=
 \norm{\E[F\circ\theta_n\mid\mathcal D']
       -\E[F\mid\mathcal D]\circ\theta_n}_2^2.
\end{align*}
The equality of norms therefore gives \eqref{eq:saturation-equality}.
\end{proof}

We fix a saturated object for the rest of the proof.  Its enlargement
need not be mixing.  All subsequent locality assertions will come from
\eqref{eq:saturation-equality} and the defining array laws. The equality
says that observing the extension's richer parameter or outside-plane
data reveals no more about an old $L^2$ test than the corresponding old
data did.

\subsection{Conditioning on parent planes}

We record the elementary conditioning fact that will justify the
induction.  It keeps track of which parameters determine each marginal.

\begin{lemma}[Local conditioning]\label{lem:local-conditioning}
Let $E,O_1,\ldots,O_s,B$ be standard random elements.  Suppose
$E_j=e_j(E)$ and $A_j=a_j(E_j,O_j)$ are measurable, and that
\[
 \Law((O_j)_{j=1}^s\mid E)
   =\bigotimes_{j=1}^s K_j(E_j,\cdot),\qquad
 \Law(B\mid E,O_1,\ldots,O_s)=L(E,A_1,\ldots,A_s;\cdot).
\]
If $Q_j$ is the conditional law of $O_j$ given $(E_j,A_j)$, then
\begin{equation}\label{eq:local-conditioning}
 \Law((O_j)_{j=1}^s\mid E,A_1,\ldots,A_s,B)
 =\bigotimes_{j=1}^s Q_j(E_j,A_j;\cdot).
\end{equation}
\end{lemma}

\begin{proof}
The tower property shows that $K_j(E_j,\cdot)$ is also the conditional
law of $O_j$ given $E_j$.  Let $M_j(x,\cdot)$ be its image under
$o\mapsto a_j(x,o)$.  Disintegration gives, for almost every $x$ with
respect to the law of $E_j$,
\[
 K_j(x,\dd o)\delta_{a_j(x,o)}(\dd\alpha)
 =M_j(x,\dd\alpha)Q_j(x,\alpha;\dd o).
\]
Substituting $x=e_j(E)$ preserves these almost-sure identities.  The
joint law of $(E,O,A,B)$ thus factors as
\[
 \Law(E)(\dd e)
 \left[\bigotimes_{j=1}^sM_j(e_j(e),\dd\alpha_j)\right]
 L(e,\alpha;\dd b)
 \left[\bigotimes_{j=1}^sQ_j(e_j(e),\alpha_j;\dd o_j)\right].
\]
The first three factors give the law of $(E,A,B)$, proving
\eqref{eq:local-conditioning}.
\end{proof}

\begin{lemma}[Rolling an axis into the local state]\label{lem:rolled-extension}
Let $R\subset I$ satisfy $|R|\leq k-2$, and choose $a\in I\setminus R$.
There is an extension of the saturated object whose local state is
$Y'=Y^I$, with local law $m_1$, parameter
\begin{equation}\label{eq:rolled-parameter}
 u'((y_c)_{c\in I})=((u(y_c))_{c\in I},(y_c)_{c\in R}),
\end{equation}
and factor map $(y_c)_{c\in I}\mapsto y_a$.
Its level-$n$ law is $m_{n+1}$, with the first axis treated as internal
to the local state, and its time transformation is componentwise $T_Y$.
\end{lemma}

\begin{proof}
We prove the statement by induction on $|R|$, simultaneously for all
array levels and all projection slots $a\notin R$. Saturation will be
used only for extensions already supplied by smaller parent sets.
The resulting screening identity and Lemma~\ref{lem:local-conditioning}
will then verify the conditional product law for the larger parent set.
Write an entry of
the old level $n+1$ as $y_{c,w}$, where $c\in I$ is the internal axis
and $w\in I^n$ comprises the remaining axes.  The proposed projection
retains $y_{a,w}$.  It preserves every array law by axis consistency,
preserves the $X$-state through $\kappa(y_a)$, and recovers the old
parameter from \eqref{eq:rolled-parameter}.  Time invariance,
consistency, and unconditional last-axis independence for the candidate
follow directly from the old level $n+1$.

At the local level, condition on $(u(y_c))_{c\in I}$ in $m_1$.  The set
$R\cup\{a\}$ is proper.  Its states are conditionally independent,
with the law at $a$ conditioned only on $u(y_a)$.  Since the $X$-image
of that conditional law is $\mu_X$, the selected $X$-state still has
law $\mu_X$ after conditioning further on $(y_c)_{c\in R}$.
This verifies the required independence from the enlarged parameter.
If $R=\varnothing$, the remaining conditional last-axis requirement is
exactly the old condition at level $n+1$, so the base case is proved.

Now suppose $R\ne\varnothing$ and that the lemma holds for smaller
parent sets.  Fix $n\geq1$ and a missing last-axis slot $j_0$.  In the
old law $m_{n+1}$ let
\[
 E=(u(y_{c,w}))_{c,w},\qquad
 O=(y_{c,w}:w_n\ne j_0).
\]
For $j\ne j_0$, let $O_j$ be the full plane at last slot $j$, $E_j$
its old parameter array, and $A_j$ its states at internal indices
$c\in R$.  Finally, let
\[
 B=(y_{c,w}:c\in R,\ w_n=j_0)
\]
be the parent states in the missing plane.  Conditional on $E$, the
$O_j$ are independent, with each plane law depending only on $E_j$.
We claim that
\begin{equation}\label{eq:parent-screening}
 \Law(B\mid E,O)
 \text{ depends on }O\text{ only through }(A_j)_{j\ne j_0}.
\end{equation}

Enumerate $R$ and reveal its missing-plane slices one internal index
$c$ at a time.  Let $R'$ be the set of indices already revealed.
It has smaller cardinality than $R$.  The induction hypothesis supplies
a valid extension with parent set $R'$ and projection slot $c$.
Apply saturation to a bounded test $F$ of the missing-plane slice
$(y_{c,w}:w_n=j_0)$, viewed as a function on the projected old level $n$.
For the data type $\mathcal G_{n,j_0}$, the old data are the parameters
of the row at $c$ and its states outside the missing plane; the enlarged
data are $E$, all of $O$, and the already revealed missing parent
slices. More explicitly, if $B_{R'}$ denotes those already revealed
missing slices, write $A_j^{R'}$ for the states of $O_j$ at internal
indices in $R'$.  This smaller-parent extension has
$\mathcal P'_n=\sigma(E,(A_j^{R'})_{j\ne j_0},B_{R'})$,
$\mathcal O'_{n,j_0}=\sigma(O)$, and
$\mathcal G'_{n,j_0}=\sigma(E,O,B_{R'})$. Hence
\eqref{eq:saturation-equality} says exactly that
\begin{equation}\label{eq:parent-reveal}
 \E[F\mid E,O,(y_{l,w})_{l\in R',w\in I^n}]
 =\E[F\mid (u(y_{c,w}))_{w\in I^n},
                (y_{c,w}:w_n\ne j_0)].
\end{equation}
The right side is the old $m_n$ conditional kernel at the fixed
internal index $c$.  It involves only $E$ and the part of $O$ belonging
to that same $c$, and it does not depend on previously revealed missing
states.  Multiplication of these successive kernels proves
\eqref{eq:parent-screening}.

Apply Lemma~\ref{lem:local-conditioning} to the proper family of planes
$O_j$.  Conditioning on $E$, all the $A_j$, and $B$ preserves their
independence and conditions each plane only on $(E_j,A_j)$.  These are
precisely the complete enlarged parameters and each plane's own
enlarged parameters, respectively.  The law of such a plane is the
old $m_n$, now viewed as the candidate level $n-1$; its conditional law
given $(E_j,A_j)$ is therefore exactly the one required by
Definition~\ref{def:array-object}.  We have proved the requirement for
$k-1$ planes.  Marginalization proves it for every proper family.

There are only countably many levels and finitely many choices of
parents and slots at each induction step.  Testing on countable
determining classes makes all the kernel equalities simultaneous.
The induction is complete.
\end{proof}

\subsection{The square and its grouped kernels}

Introduce the disintegrations
\begin{equation}\label{eq:local-and-row-disintegrations}
 m_0(\dd y)=\int_U\eta_u(\dd y)\nu(\dd u),\qquad
 m_1(\dd\mathbf y)=\int_{U^I}\lambda_p(\dd\mathbf y)\pi(\dd p),
\end{equation}
where $\pi$ is the law of the parameter tuple under $m_1$.
The local law $\eta_u$ is concentrated on $u(y)=u$ and satisfies
\begin{equation}\label{eq:fibre-X-marginal}
 \kappa_*\eta_u=\mu_X\quad\text{for }\nu\text{-almost every }u.
\end{equation}
Every proper coordinate marginal of $\pi$ is a product of copies of
$\nu$.  Every proper coordinate marginal of $\lambda_p$ is
$\bigotimes_j\eta_{p_j}$ on those coordinates.  These assertions follow
respectively from unconditional and conditional last-axis independence
at level one.

The measure $\nu$ is $T_U$-invariant, and $\pi$ is invariant under
$T_U^I$.  Uniqueness of disintegration and time invariance give
\begin{equation}\label{eq:fibre-time-equivariance}
 \eta_{T_Uu}=(T_Y)_*\eta_u,\qquad
 \lambda_{T_U^Ip}=(T_Y^I)_*\lambda_p.
\end{equation}
By intersecting the relevant conull sets over all integer translates,
we may use these identities, the marginal assertions, and
\eqref{eq:fibre-X-marginal} on invariant conull sets.  In particular,
time carries one fibre Hilbert space to another; it need not preserve
an individual fibre.

Write the level-two state as $(y_{a,j})_{a,j\in I}$, with $a$ the first
axis (the row index) and $j$ the last axis (the column index).  Its full
parameter array and its row and column parameter tuples are
\[
 E=(u_{a,j})_{a,j\in I},\qquad
 p_a=(u_{a,j})_{j\in I},\qquad q_j=(u_{a,j})_{a\in I}.
\]
Thus the order within either tuple is the original order on $I$.
For subsets $D,J\subset I$, let
$Y_{D,J}=(y_{a,j})_{a\in D,j\in J}$; write $Y_{a,I}$ for a single row
and $Y_{I,j}$ for a single column.

\begin{proposition}[Square locality]\label{prop:square-locality}
The following properties hold in $m_2$, and can be imposed
simultaneously on one conull set of full parameter arrays $E$.
\begin{enumerate}
\item Conditional on $E$, every proper family of columns is independent,
with column $j$ having law $\lambda_{q_j}$.  Every proper family of rows
is independent, with row $a$ having law $\lambda_{p_a}$.
\item Let $R\subset I$, $|R|\leq k-2$.  Given $E$ and the full parent
rows $Y_{R,I}$, every proper family of columns is independent, and each
column has law $\lambda_{q_j}$ conditioned on its own parent entries
$Y_{R,j}$.
\item If $a\notin R$ and $j_0\in I$, then
\begin{equation}\label{eq:single-cell-locality}
 \Law(y_{a,j_0}\mid E,Y_{R,I},Y_{I,I\setminus\{j_0\}})
 =\lambda_{p_a}(\dd y_{a,j_0}\mid Y_{a,I\setminus\{j_0\}}).
\end{equation}
\end{enumerate}
There is also the following grouped form.  Partition the rows into
three nonempty groups $P,A,B$ and the columns into three nonempty groups
$r,s,t$.  At a fixed typical $E$, each grouped row
\[
 (Y_{D,r},Y_{D,s},Y_{D,t}),\quad D\in\{P,A,B\},
\]
has a three-coordinate law with pairwise independent coordinates, and
the same holds for each grouped column.  Given $Y_{P,I}$, any two
grouped columns are independent with their own column laws conditioned
on their entries in $P$.  Finally, putting
$\Lambda_A=\bigotimes_{a\in A}\lambda_{p_a}$, we have
\begin{equation}\label{eq:grouped-cell-locality}
 \Law(Y_{A,t}\mid E,Y_{P,I},Y_{I,r\cup s})
 =\Lambda_A(\dd Y_{A,t}\mid Y_{A,r},Y_{A,s}).
\end{equation}
All these grouped assertions hold for every ordering of the three row
groups and the three column groups.
\end{proposition}

\begin{proof}
Column independence in (1) is the defining conditional independence at
level two.  To obtain rows, use the extension in
Lemma~\ref{lem:rolled-extension} with parent set $R$ and projection
slot $a\notin R$.  For a bounded test $F$ of the entire projected row,
saturation for the full parameter data $\mathcal P_1$ gives
\begin{equation}\label{eq:row-given-parents}
 \E[F(Y_{a,I})\mid E,Y_{R,I}]
 =\int F\dd\lambda_{p_a}.
\end{equation}
Reveal the members of any proper row family in turn.  The previous
members number at most $k-2$, so \eqref{eq:row-given-parents} applies
at every step.  Their joint conditional law is the product of the
indicated row laws.

Assertion (2) is the conditional last-axis requirement at level one of
the extension with parents $R$.  Its local law is $m_1$ on each column,
and its own parameter there is $(q_j,Y_{R,j})$, giving precisely the
stated marginal kernel.  For (3), use the same extension and saturation
for $\mathcal G_{1,j_0}$.  The projected old data are $p_a$ and the
other entries of row $a$, while the extension data are $E$, the parents,
and all columns except $j_0$.  This gives
\eqref{eq:single-cell-locality}.

We prove the grouped assertions without interchanging array axes.
Each row group is a proper set of rows, so its full law is the product
of its individual row laws.  Within each individual row, any proper
set of column coordinates is independent.  Consequently any two
grouped cells in that grouped row are independent.  The identical
conclusion for a grouped column follows by applying the corresponding
column assertions.  Since $|P|\leq k-2$, (2) applies to the parents
$P$.  The union of any two column groups is proper, which gives the
claimed conditional independence of grouped columns.

It remains to prove \eqref{eq:grouped-cell-locality}.  We first handle
a single missing column $j_0$.  Given $E$, $Y_{P,I}$, and all other
columns, reveal $y_{a,j_0}$ successively for $a\in A$.  At each step
the parents consist of $P$ and the previously revealed rows in $A$.
Their number is at most
\[
 |P|+|A|-1=k-|B|-1\leq k-2.
\]
All their remaining entries are already known from the observed
columns.  Equation~\eqref{eq:single-cell-locality} therefore gives
the kernel of the individual row at that step.  These kernels do not
depend on the previously revealed entries in column $j_0$, and hence
\begin{equation}\label{eq:group-single-missing}
 \Law(Y_{A,j_0}\mid E,Y_{P,I},Y_{I,I\setminus\{j_0\}})
 =\bigotimes_{a\in A}
   \lambda_{p_a}(\dd y_{a,j_0}\mid Y_{a,I\setminus\{j_0\}}).
\end{equation}

Now choose $j_0\in t$, put $J=t\setminus\{j_0\}$ and
$C=r\cup s$, and retain
$W=Y_{A,J}$ as part of the output.  The set $C\cup J$ consists of
exactly $k-1$ columns.  By (2), conditional on $E$ and $Y_{P,I}$,
those columns are independent.  In each of them, $P\cup A$ is a
proper row set.  Thus the $A$-entries, even after conditioning on its
$P$-entries, have the plain product marginal
$\bigotimes_{a\in A}\eta_{u_{a,j}}$.  It follows that
\begin{equation}\label{eq:group-intermediate-marginal}
 \Law(W\mid E,Y_{P,I},Y_{I,C})
 =\bigotimes_{a\in A,\,j\in J}\eta_{u_{a,j}}.
\end{equation}
Use \eqref{eq:group-single-missing} and then integrate away the
additional states in columns $J$, retaining $W$.  Its right side
depends on those columns only through $W$, so the resulting joint
kernel for $(W,Y_{A,j_0})$ is the measure in
\eqref{eq:group-intermediate-marginal}, followed by the product of the
individual row kernels in \eqref{eq:group-single-missing}.

Under the row-group law $\Lambda_A$, the coordinates in $C\cup J$
are also independent with exactly those marginals: the rows in $A$
are independent and each row has independent proper coordinates.
Consequently the just obtained kernel is exactly the conditional law
of $Y_{A,t}$ under $\Lambda_A$ given $Y_{A,C}$.  This proves
\eqref{eq:grouped-cell-locality}, including the case $J=\varnothing$.
Only cardinalities of the groups were used, so every ordering of the
groups is allowed.  All choices are finite; the earlier convention on
kernel versions makes the assertions simultaneous.
\end{proof}

\subsection{The geometry of the parameters}

The last property needed from the square is unconditional: it describes
how a row parameter tuple and the column parameter tuples can meet.
For $a\in I$, let
\begin{equation}\label{eq:parameter-coordinate-kernel}
 K_a(u,\dd q)=\pi(\dd q\mid q[a]=u).
\end{equation}
This is a probability kernel for $\nu$-almost every $u$, because the
$a$th coordinate marginal of $\pi$ is $\nu$.

\begin{lemma}[Parameter geometry]\label{lem:parameter-geometry}
In the level-two parameter law, for every $a\in I$ and every proper
set $J\subset I$,
\begin{equation}\label{eq:parameter-geometry}
 \Law((q_j)_{j\in J}\mid p_a)
 =\bigotimes_{j\in J}K_a(p_a[j],\dd q_j).
\end{equation}
In particular, for each $i\in I$, the diagonal row and column tuples
$p_i,q_i$ are conditionally independent, with the same law
$K_i(u,\cdot)$, given their common coordinate $u=u_{i,i}$.
\end{lemma}

\begin{proof}
It is enough to prove \eqref{eq:parameter-geometry} for
$J=I\setminus\{j_0\}$, and then take marginals.  Write
\[
 Z=Y_{a,I},\qquad O=(Y_{I,j})_{j\in J},\qquad
 S=(y_{a,j})_{j\in J}.
\]
Use the extension with no parents and projection slot $a$.
Saturation for the data $\mathcal O_{1,j_0}$ says, for every bounded
row test $F$,
\begin{equation}\label{eq:unconditional-row-screening}
 \E[F(Z)\mid O]=\E[F(Z)\mid S].
\end{equation}
Since $S$ is a function of both $Z$ and $O$, this is conditional
independence of $Z$ and $O$ given $S$.  One direct verification is to
multiply \eqref{eq:unconditional-row-screening} by a bounded function
of $O$ and condition on $S$; the resulting factorization is the
conditional-independence identity.  It also gives the reversed
screening statement $\Law(O\mid Z)=\Law(O\mid S)$.

Unconditionally the columns in $O$ are independent, each with law
$m_1$.  Conditioning on their separate shared states therefore gives
\[
 \Law(O\mid Z)
 =\bigotimes_{j\in J}m_1(\dd Y_{I,j}\mid y_{a,j}).
\]
We identify the parameter image of each factor.  In a single column,
the joint law of its parameter tuple $q$ and its state $y_a$ at slot
$a$ is
\begin{equation}\label{eq:shared-state-law}
 \pi(\dd q)\eta_{q[a]}(\dd y_a).
\end{equation}
Disintegrating $\pi$ at $q[a]$ rewrites this as
\[
 \nu(\dd u)K_a(u,\dd q)\eta_u(\dd y_a).
\]
Since $u(y_a)=u$ almost surely under $\eta_u$, the conditional law
of $q$ given the shared state $y_a$ is exactly
$K_a(u(y_a),\dd q)$.  Thus
\[
 \Law((q_j)_{j\in J}\mid Z)
 =\bigotimes_{j\in J}K_a(u(y_{a,j}),\dd q_j).
\]
The right side is a function only of $p_a$, which itself is a
function of $Z$.  Conditioning down to $p_a$ proves
\eqref{eq:parameter-geometry}.

For the last assertion, take $a=i$ and retain just $j=i$.  The tuple
$p_i$ has marginal $\pi$, and its coordinate at $i$ is $u_{i,i}$.
Given $p_i$, the law of $q_i$ is $K_i(p_i[i],\cdot)$.
Disintegrating once more over that coordinate shows that
\[
 \Law(p_i,q_i\mid u_{i,i}=u)
 =K_i(u,\dd p_i)K_i(u,\dd q_i),
\]
as required.
\end{proof}

We now have two complementary descriptions of the same square.
Proposition~\ref{prop:square-locality} gives exact conditional kernels
at a fixed full parameter array.  Lemma~\ref{lem:parameter-geometry}
controls how those parameters vary.  The first description yields the
operator structure in the next section; the second will identify which
parts of that structure persist across fibres.
